\documentclass[11pt]{article}
\usepackage{mathblatt}
% gesetzt von werkzeuge/setzer.py (Sorte selbst) aus dem Lernweg HDK
% und den Bankzeilen; Muster bau/proben/2026-10-09/hypotenuse-muster.
\usepackage{enumitem}
\usepackage{adjustbox,varwidth}
\newgeometry{a4paper,margin=18mm,top=15mm,bottom=20mm,footskip=9mm}
\pagestyle{fancy}\fancyhf{}\renewcommand{\headrulewidth}{0pt}
\fancyfoot[L]{\footnotesize\color{mbgrau}HDK \textperiodcentered{} Winkelsummen, Dreiecke und Vierecke}
\fancyfoot[R]{\footnotesize\color{mbgrau}\thepage}
\setlength{\parskip}{3pt}
\renewcommand{\kreuz}[1]{\mbox{$\square$\ #1}\quad}
\newcommand{\kreuzl}[1]{\par\hangindent1.4em\hangafter1\noindent$\square$\ #1}
\newcommand{\aufg}[3]{\par\noindent\makebox[0pt][r]{#3}\begin{minipage}[t]{\linewidth}#1\end{minipage}\par}
\newcommand{\aufgg}[4]{\par\noindent\makebox[0pt][r]{#4}\begin{minipage}[t]{0.6\linewidth}\vspace{0pt}#1\end{minipage}\hfill
  \begin{minipage}[t]{0.37\linewidth}\vspace{0pt}\centering #2\end{minipage}\par}
\newcommand{\nr}[1]{\makebox[7mm][l]{\textbf{#1}}}
\newcommand{\lsg}[3]{\par\noindent\begin{minipage}[t]{9mm}\textbf{#1}\end{minipage}%
  \begin{minipage}[t]{52mm}\raggedright\bfseries\boldmath #2\end{minipage}\hspace{3mm}%
  \begin{minipage}[t]{\dimexpr\linewidth-64mm\relax}\raggedright\small #3\end{minipage}\par\vspace{4pt}}
\newlength{\kr}
\makeatletter
\newcommand{\karomerk}[2]{\protected@write\@auxout{}{\string\karoseite{#1}{#2}{\thepage}}}
\newcommand{\karoseite}[3]{\expandafter\xdef\csname karo@#1@#2\endcsname{#3}}
\newcommand{\karohinweis}[3]{\karomerk{h}{#1}\ifcsname karo@k@#1\endcsname
  \ifnum\csname karo@k@#1\endcsname=0\csname karo@h@#1\endcsname\relax#2\else#3\fi
  \else#3\fi}
\makeatother
\begin{document}

\noindent{\LARGE\bfseries Winkelsummen, Dreiecke und Vierecke}\par\smallskip
\noindent{\small\color{mbgrau}Selbstlernheft Klasse 7 \textperiodcentered{} mit Lösungsheft \textperiodcentered{} Kennung HDK}\par\medskip
\noindent\textbf{So nutzt du das Heft:} Lies im Abschnitt das Beispiel. Rechne dann die Aufgaben der Reihe nach und vergleiche jede sofort mit dem Lösungsheft. Kreuze „kann ich“ an, wenn du die letzte Aufgabe (\emph{Ziel}) allein schaffst.\par
\noindent Mehr Aufgaben zu einem Abschnitt: Kennung deinem Nachhilfelehrer schicken (z.\,B. „HDK mehr B“).\par\medskip
{\arrayrulecolor{mbgrau}\renewcommand{\arraystretch}{1.3}
\noindent\begin{tabular}{|>{\bfseries}p{10mm}|p{112mm}|>{\raggedleft\arraybackslash}p{10mm}|>{\centering\arraybackslash}p{14mm}|}\hline
\rowcolor{black!8} & \textbf{Abschnitt} & \textbf{Seite} & \textbf{kann ich}\\\hline
A & Winkelsumme im Dreieck & \pageref{s:A} & $\square$\\\hline
B & Gleichschenklige Dreiecke & \pageref{s:B} & $\square$\\\hline
C & Winkelsumme im Viereck & \pageref{s:C} & $\square$\\\hline
D & Winkel in Teildreiecken & \pageref{s:D} & $\square$\\\hline
T & Probetest & \pageref{s:T} & $\square$\\\hline
\end{tabular}}\par\bigskip

\begin{tcolorbox}[colback=mbkasten,colframe=mbgrau,boxrule=0.5pt,arc=1mm,left=3mm,right=3mm,top=2mm,bottom=2mm,title={\bfseries Formel auf einen Blick},coltitle=black,colbacktitle=black!10]
\begin{minipage}[t]{0.58\linewidth}\vspace{0pt}
{\Large Dreieck: zusammen $180^\circ$\par Viereck: zusammen $360^\circ$}\par\medskip
In Worten: Fehlt ein Winkel, ziehst du die bekannten von $180^\circ$ (Dreieck) oder $360^\circ$ (Viereck) ab. Gleich lange Seiten oder eine Symmetrieachse sagen dir, welche Winkel gleich sind.\par\medskip
\textbf{So gehst du vor}\par
\nr{1}das Dreieck oder Viereck suchen, in dem der gesuchte Winkel liegt\par
\nr{2}bekannte Winkel sammeln, auch gleiche Winkel und Nebenwinkel\par
\nr{3}von $180^\circ$ oder $360^\circ$ abziehen\par
\nr{4}Probe: alle Winkel zusammen\par
\medskip
\end{minipage}\hfill
\begin{minipage}[t]{0.38\linewidth}\vspace{2mm}\centering
\begin{tikzpicture}[line width=0.7pt,font=\footnotesize]\draw (0.000,0.000) -- (3.600,0.000) -- (2.489,2.089) -- cycle;\draw (2.489,2.089) -- (1.930,0.000);\draw[mbgrau,line width=0.5pt] (0.380,0.000) arc (0.00:40.00:0.38);\node[inner sep=0pt] at (0.788,0.287) {$40^\circ$};\draw[mbgrau,line width=0.5pt] (3.422,0.336) arc (118.00:180.00:0.38);\node[inner sep=0pt] at (2.980,0.373) {$62^\circ$};\draw[mbgrau,line width=0.5pt] (2.198,1.845) arc (-140.00:-105.00:0.38);\node[inner sep=0pt] at (1.922,1.199) {$35^\circ$};\draw[mbgrau,line width=0.5pt] (2.391,1.722) arc (-105.00:-62.00:0.38);\node[inner sep=0pt] at (2.558,1.488) {$\varepsilon$};\node[inner sep=0pt] at (-0.235,-0.086) {$A$};\node[inner sep=0pt] at (3.814,-0.129) {$B$};\node[inner sep=0pt] at (2.537,2.334) {$C$};\node[inner sep=0pt] at (1.930,-0.220) {$D$};\end{tikzpicture}
\end{minipage}
\end{tcolorbox}
\medskip\noindent\textbf{Typische Fehler – so nicht}\par\smallskip
\noindent\nr{$\times$}Im Viereck mit $180^\circ$ gerechnet. Ein Viereck hat zusammen $360^\circ$.\par
\noindent\nr{$\times$}Im gleichschenkligen Dreieck den Winkel an der Spitze so groß wie einen Basiswinkel gesetzt ($68^\circ$ statt $44^\circ$).\par
\noindent\nr{$\times$}Im Teildreieck mit einem Winkel des großen Dreiecks gerechnet. Erst das Teildreieck nachfahren, dann rechnen.\par
\clearpage\label{s:A}\noindent\colorbox{black!12}{\makebox[10mm]{\Large\bfseries\strut A}}\hspace{3mm}{\Large\bfseries Winkelsumme im Dreieck}\par\smallskip
\noindent Die drei Winkel eines Dreiecks ergeben zusammen $180^\circ$, egal wie das Dreieck aussieht.\par\medskip
\noindent\textbf{Formel}\par\nopagebreak\noindent\hspace*{4mm}\begin{minipage}{\dimexpr\linewidth-4mm\relax}$\alpha + \beta + \gamma = 180^\circ$ \qquad $\gamma = 180^\circ - \alpha - \beta$\par {\small Ein Bogen mit Punkt ist ein rechter Winkel: $90^\circ$.}\end{minipage}\par\medskip
\noindent\textbf{Beispiel}\quad (a) Im Dreieck $ABC$ ist $\alpha = 52{,}4^\circ$ und $\beta = 71{,}3^\circ$. Berechne $\gamma$. (b) Im Dreieck $DEF$ ist der Winkel bei $F$ ein rechter, der Winkel bei $D$ ist $34^\circ$. Berechne $\varepsilon$.\par\smallskip
\noindent\begin{minipage}[t]{0.6\linewidth}\vspace{0pt}\par\noindent\hangindent6mm\hangafter1\makebox[6mm][l]{\textbf{1}}(a) Die bekannten Winkel zusammen: \ $52{,}4^\circ + 71{,}3^\circ = 123{,}7^\circ$\par\smallskip\par\noindent\hangindent6mm\hangafter1\makebox[6mm][l]{\textbf{2}}Der Rest bis $180^\circ$ ist $\gamma$: \ $\gamma = 180^\circ - 123{,}7^\circ = 56{,}3^\circ$\par\smallskip\par\noindent\hangindent6mm\hangafter1\makebox[6mm][l]{\textbf{3}}(b) Bogen mit Punkt bei $F$: ein rechter Winkel, $90^\circ$. \ $90^\circ + 34^\circ = 124^\circ$\par\smallskip\par\noindent\hangindent6mm\hangafter1\makebox[6mm][l]{\textbf{4}}$\varepsilon = 180^\circ - 124^\circ = 56^\circ$\par\smallskip\end{minipage}\hfill
\begin{minipage}[t]{0.37\linewidth}\vspace{0pt}\centering \begin{tikzpicture}[line width=0.7pt,font=\footnotesize]\draw (0.000,0.000) -- (2.882,0.000) -- (2.002,2.600) -- cycle;\draw[mbgrau,line width=0.5pt] (0.380,0.000) arc (0.00:52.40:0.38);\node[inner sep=0pt] at (0.831,0.409) {$52{,}4^\circ$};\draw[mbgrau,line width=0.5pt] (2.760,0.360) arc (108.70:180.00:0.38);\node[inner sep=0pt] at (2.212,0.481) {$71{,}3^\circ$};\draw[mbgrau,line width=0.5pt] (1.770,2.299) arc (-127.60:-71.30:0.38);\node[inner sep=0pt] at (1.902,2.000) {$\gamma$};\node[inner sep=0pt] at (-0.224,-0.110) {$A$};\node[inner sep=0pt] at (3.085,-0.146) {$B$};\node[inner sep=0pt] at (2.043,2.847) {$C$};\draw (4.300,0.000) -- (7.100,0.000) -- (6.224,1.298) -- cycle;\draw[mbgrau,line width=0.5pt] (4.680,0.000) arc (0.00:34.00:0.38);\node[inner sep=0pt] at (5.214,0.279) {$34^\circ$};\draw[mbgrau,line width=0.5pt] (5.976,1.130) arc (-146.00:-56.00:0.3);\fill (6.199,1.166) circle (0.8pt);\draw[mbgrau,line width=0.5pt] (6.888,0.315) arc (124.00:180.00:0.38);\node[inner sep=0pt] at (6.568,0.283) {$\varepsilon$};\node[inner sep=0pt] at (4.061,-0.073) {$D$};\node[inner sep=0pt] at (7.321,-0.117) {$E$};\node[inner sep=0pt] at (6.272,1.543) {$F$};\end{tikzpicture}\end{minipage}\par\medskip
\noindent\textbf{Aufgaben}\hspace{1em}{\small\color{mbgrau}\karohinweis{A}{Rechne unten im Karo.}{Rechne auf einem extra Blatt.} Vergleiche jede Aufgabe gleich mit dem Lösungsheft.}\par\smallskip
\aufgg{\hangindent7mm\hangafter1\nr{1.}Berechne den Winkel $\gamma$.}{\begin{tikzpicture}[line width=0.7pt,font=\footnotesize]\draw (0.000,0.000) -- (2.743,0.000) -- (2.944,2.300) -- cycle;\draw[mbgrau,line width=0.5pt] (0.380,0.000) arc (0.00:38.00:0.38);\node[inner sep=0pt] at (0.821,0.283) {$38^\circ$};\draw[mbgrau,line width=0.5pt] (2.776,0.379) arc (85.00:180.00:0.38);\node[inner sep=0pt] at (2.261,0.525) {$95^\circ$};\draw[mbgrau,line width=0.5pt] (2.644,2.066) arc (-142.00:-95.00:0.38);\node[inner sep=0pt] at (2.647,1.754) {$\gamma$};\node[inner sep=0pt] at (-0.236,-0.081) {$A$};\node[inner sep=0pt] at (2.912,-0.184) {$B$};\node[inner sep=0pt] at (3.063,2.520) {$C$};\end{tikzpicture}}{}{}
\medskip
\aufgg{\hangindent7mm\hangafter1\nr{2.}Berechne den Winkel $\beta$.}{\begin{tikzpicture}[line width=0.7pt,font=\footnotesize]\draw (0.000,0.000) -- (2.509,0.000) -- (1.346,2.700) -- cycle;\draw[mbgrau,line width=0.5pt] (0.380,0.000) arc (0.00:63.50:0.38);\node[inner sep=0pt] at (0.706,0.437) {$63{,}5^\circ$};\draw[mbgrau,line width=0.5pt] (2.359,0.349) arc (113.30:180.00:0.38);\node[inner sep=0pt] at (2.000,0.335) {$\beta$};\draw[mbgrau,line width=0.5pt] (1.177,2.360) arc (-116.50:-66.70:0.38);\node[inner sep=0pt] at (1.317,1.646) {$49{,}8^\circ$};\node[inner sep=0pt] at (-0.213,-0.132) {$A$};\node[inner sep=0pt] at (2.718,-0.137) {$B$};\node[inner sep=0pt] at (1.353,2.950) {$C$};\end{tikzpicture}}{}{}
\medskip
\aufgg{\hangindent7mm\hangafter1\nr{3.}Das Dreieck $ABC$ hat bei $B$ einen rechten Winkel. Berechne $\gamma$.}{\begin{tikzpicture}[line width=0.7pt,font=\footnotesize]\draw (0.000,0.000) -- (3.000,0.000) -- (3.000,1.562) -- cycle;\draw[mbgrau,line width=0.5pt] (0.380,0.000) arc (0.00:27.50:0.38);\node[inner sep=0pt] at (1.335,0.327) {$27{,}5^\circ$};\draw[mbgrau,line width=0.5pt] (3.000,0.300) arc (90.00:180.00:0.3);\fill (2.905,0.095) circle (0.8pt);\draw[mbgrau,line width=0.5pt] (2.663,1.386) arc (-152.50:-90.00:0.38);\node[inner sep=0pt] at (2.677,1.030) {$\gamma$};\node[inner sep=0pt] at (-0.243,-0.059) {$A$};\node[inner sep=0pt] at (3.177,-0.177) {$B$};\node[inner sep=0pt] at (3.130,1.775) {$C$};\end{tikzpicture}}{}{}
\medskip
\aufg{\hangindent7mm\hangafter1\nr{4.}Ein dreieckiges Sonnensegel wird an drei Punkten gespannt. An zwei Ecken misst Jonas $66{,}5^\circ$ und $59{,}5^\circ$. Der Hersteller schreibt: Keine Ecke darf spitzer als $55^\circ$ sein, sonst reißt die Naht. Hält das Segel? Rechne nach.}{}{{\scriptsize\color{mbgrau}Ziel\hspace{2mm}}}
\medskip
\par\vspace{3mm}\setlength{\kr}{\dimexpr\pagegoal-\pagetotal-10mm\relax}%
\ifdim\kr>12mm\karomerk{k}{A}\noindent{\scriptsize\color{mbgrau}Platz zum Rechnen}\par\nointerlineskip\vspace{1mm}\noindent
\begin{tikzpicture}\clip (0,0) rectangle ({\linewidth-0.5mm},\kr);\draw[black!16,line width=0.3pt,step=5mm] (0,0) grid (\linewidth,\kr);\end{tikzpicture}\fi\par
\clearpage\label{s:B}\noindent\colorbox{black!12}{\makebox[10mm]{\Large\bfseries\strut B}}\hspace{3mm}{\Large\bfseries Gleichschenklige Dreiecke}\par\smallskip
\noindent Im \textbf{gleichschenkligen} Dreieck sind zwei Seiten gleich lang (Striche). Die beiden Winkel an der dritten Seite, die \textbf{Basiswinkel}, sind gleich groß. Umgekehrt: Sind zwei Winkel gleich, sind die Seiten gegenüber gleich lang.\par\medskip
\noindent\textbf{Formel}\par\nopagebreak\noindent\hspace*{4mm}\begin{minipage}{\dimexpr\linewidth-4mm\relax}Spitze $= 180^\circ - 2 \cdot$ Basiswinkel \qquad Basiswinkel $= (180^\circ - \text{Spitze}) : 2$\par {\small Gleichseitig (drei gleiche Seiten): jeder Winkel $60^\circ$.}\end{minipage}\par\medskip
\noindent\textbf{Beispiel}\quad Im Dreieck $ABC$ sind $AC$ und $BC$ gleich lang. Es ist $\alpha = 68^\circ$. Berechne $\beta$ und $\gamma$.\par\smallskip
\noindent\begin{minipage}[t]{0.6\linewidth}\vspace{0pt}\par\noindent\hangindent6mm\hangafter1\makebox[6mm][l]{\textbf{1}}$AC$ und $BC$ gleich lang (Striche): Die Winkel an der dritten Seite $AB$ sind die \textbf{Basiswinkel}. \ $\beta = \alpha = 68^\circ$\par\smallskip\par\noindent\hangindent6mm\hangafter1\makebox[6mm][l]{\textbf{2}}$\gamma = 180^\circ - 68^\circ - 68^\circ = 44^\circ$\par\smallskip\par\noindent\hangindent6mm\hangafter1\makebox[6mm][l]{\textbf{3}}Ist die Spitze gegeben, z.\,B. $\gamma = 52^\circ$: \ $(180^\circ - 52^\circ) : 2 = 64^\circ$, also $\alpha = \beta = 64^\circ$\par\smallskip\end{minipage}\hfill
\begin{minipage}[t]{0.37\linewidth}\vspace{0pt}\centering \begin{tikzpicture}[line width=0.7pt,font=\footnotesize]\draw (0.000,0.000) -- (1.859,0.000) -- (0.929,2.300) -- cycle;\draw[line width=0.6pt] (0.372,1.187) -- (0.557,1.113);\draw[line width=0.6pt] (1.301,1.113) -- (1.487,1.187);\draw[mbgrau,line width=0.5pt] (0.380,0.000) arc (0.00:68.00:0.38);\node[inner sep=0pt] at (0.599,0.404) {$68^\circ$};\draw[mbgrau,line width=0.5pt] (1.716,0.352) arc (112.00:180.00:0.38);\node[inner sep=0pt] at (1.353,0.341) {$\beta$};\draw[mbgrau,line width=0.5pt] (0.787,1.948) arc (-112.00:-68.00:0.38);\node[inner sep=0pt] at (0.929,1.705) {$\gamma$};\node[inner sep=0pt] at (-0.207,-0.140) {$A$};\node[inner sep=0pt] at (2.066,-0.140) {$B$};\node[inner sep=0pt] at (0.929,2.550) {$C$};\end{tikzpicture}\end{minipage}\par\medskip
\noindent\textbf{Aufgaben}\hspace{1em}{\small\color{mbgrau}\karohinweis{B}{Rechne unten im Karo.}{Rechne auf einem extra Blatt.} Vergleiche jede Aufgabe gleich mit dem Lösungsheft.}\par\smallskip
\aufgg{\hangindent7mm\hangafter1\nr{1.}Im Dreieck $ABC$ sind $AC$ und $BC$ gleich lang. Berechne $\alpha$ und $\gamma$.}{\begin{tikzpicture}[line width=0.7pt,font=\footnotesize]\draw (0.000,0.000) -- (1.766,0.000) -- (0.883,2.300) -- cycle;\draw[line width=0.6pt] (0.348,1.186) -- (0.535,1.114);\draw[line width=0.6pt] (1.231,1.114) -- (1.418,1.186);\draw[mbgrau,line width=0.5pt] (0.380,0.000) arc (0.00:69.00:0.38);\node[inner sep=0pt] at (0.503,0.346) {$\alpha$};\draw[mbgrau,line width=0.5pt] (1.630,0.355) arc (111.00:180.00:0.38);\node[inner sep=0pt] at (1.170,0.409) {$69^\circ$};\draw[mbgrau,line width=0.5pt] (0.747,1.945) arc (-111.00:-69.00:0.38);\node[inner sep=0pt] at (0.883,1.705) {$\gamma$};\node[inner sep=0pt] at (-0.206,-0.142) {$A$};\node[inner sep=0pt] at (1.972,-0.142) {$B$};\node[inner sep=0pt] at (0.883,2.550) {$C$};\end{tikzpicture}}{}{}
\medskip
\aufgg{\hangindent7mm\hangafter1\nr{2.}Im Dreieck $ABC$ sind $AB$ und $AC$ gleich lang, $\alpha = 38^\circ$. Berechne $\beta$ und $\gamma$.}{\begin{tikzpicture}[line width=0.7pt,font=\footnotesize]\draw (2.300,-1.100) -- (2.300,1.100) -- (0.000,0.000) -- cycle;\draw[line width=0.6pt] (1.107,-0.640) -- (1.193,-0.460);\draw[line width=0.6pt] (1.193,0.460) -- (1.107,0.640);\draw[mbgrau,line width=0.5pt] (0.343,-0.164) arc (-25.56:25.56:0.38);\node[inner sep=0pt] at (0.686,0.000) {$38^\circ$};\draw[mbgrau,line width=0.5pt] (2.300,-0.720) arc (90.00:154.44:0.38);\node[inner sep=0pt] at (1.968,-0.574) {$\beta$};\draw[mbgrau,line width=0.5pt] (1.957,0.936) arc (-154.44:-90.00:0.38);\node[inner sep=0pt] at (1.968,0.574) {$\gamma$};\node[inner sep=0pt] at (2.433,-1.312) {$B$};\node[inner sep=0pt] at (2.433,1.312) {$C$};\node[inner sep=0pt] at (-0.250,0.000) {$A$};\end{tikzpicture}}{}{}
\medskip
\aufgg{\hangindent7mm\hangafter1\nr{3.}Im Dreieck $ABC$ ist $\alpha = \beta = 65^\circ$. Die Seite $AC$ ist 5\,cm lang. Welche Seite ist auch 5\,cm lang? Berechne $\gamma$.}{\begin{tikzpicture}[line width=0.7pt,font=\footnotesize]\draw (0.000,0.000) -- (2.145,0.000) -- (1.073,2.300) -- cycle;\node[inner sep=0pt] at (0.058,1.373) {5\,cm};\draw[mbgrau,line width=0.5pt] (0.380,0.000) arc (0.00:65.00:0.38);\node[inner sep=0pt] at (0.610,0.389) {$65^\circ$};\draw[mbgrau,line width=0.5pt] (1.984,0.344) arc (115.00:180.00:0.38);\node[inner sep=0pt] at (1.535,0.389) {$65^\circ$};\draw[mbgrau,line width=0.5pt] (0.912,1.956) arc (-115.00:-65.00:0.38);\node[inner sep=0pt] at (1.073,1.705) {$\gamma$};\node[inner sep=0pt] at (-0.211,-0.134) {$A$};\node[inner sep=0pt] at (2.356,-0.134) {$B$};\node[inner sep=0pt] at (1.073,2.550) {$C$};\end{tikzpicture}}{}{}
\medskip
\aufg{\hangindent7mm\hangafter1\nr{4.}Bei einer Stehleiter sind beide Holme gleich lang; unten stehen sie auf dem Boden, oben treffen sie sich unter $32^\circ$. Die Leiter steht sicher, wenn jeder Holm mit dem Boden mindestens $70^\circ$ bildet. Kreuze an.\par\smallskip\leftskip7mm \kreuzl{Sicher: Jeder Holm bildet $74^\circ$ mit dem Boden.}\kreuzl{Nicht sicher: Jeder Holm bildet nur $58^\circ$ mit dem Boden.}\kreuzl{Nicht sicher: Jeder Holm bildet nur $32^\circ$ mit dem Boden.}}{}{{\scriptsize\color{mbgrau}Ziel\hspace{2mm}}}
\medskip
\par\vspace{3mm}\setlength{\kr}{\dimexpr\pagegoal-\pagetotal-10mm\relax}%
\ifdim\kr>12mm\karomerk{k}{B}\noindent{\scriptsize\color{mbgrau}Platz zum Rechnen}\par\nointerlineskip\vspace{1mm}\noindent
\begin{tikzpicture}\clip (0,0) rectangle ({\linewidth-0.5mm},\kr);\draw[black!16,line width=0.3pt,step=5mm] (0,0) grid (\linewidth,\kr);\end{tikzpicture}\fi\par
\clearpage\label{s:C}\noindent\colorbox{black!12}{\makebox[10mm]{\Large\bfseries\strut C}}\hspace{3mm}{\Large\bfseries Winkelsumme im Viereck}\par\smallskip
\noindent Die vier Winkel eines Vierecks ergeben zusammen $360^\circ$: Eine Diagonale teilt es in zwei Dreiecke. Im \textbf{Drachen} sind die beiden Winkel neben der Symmetrieachse gleich groß.\par\medskip
\noindent\textbf{Formel}\par\nopagebreak\noindent\hspace*{4mm}\begin{minipage}{\dimexpr\linewidth-4mm\relax}$\alpha + \beta + \gamma + \delta = 360^\circ$\par {\small Trapez: ein Paar paralleler Seiten. Parallelogramm: zwei Paare. Raute: Parallelogramm mit vier gleich langen Seiten. Rechteck: vier rechte Winkel. Drachen: zwei Paare gleich langer Nachbarseiten, Diagonalen senkrecht.}\end{minipage}\par\medskip
\noindent\textbf{Beispiel}\quad (a) Im Viereck $ABCD$ ist $\alpha = 85^\circ$, $\beta = 110^\circ$, $\gamma = 72^\circ$. Berechne $\delta$. (b) $KLMN$ ist ein Drachen mit der Symmetrieachse $KM$. Bei $K$ sind es $64^\circ$, bei $M$ $100^\circ$. Berechne die Winkel bei $L$ und $N$.\par\smallskip
\noindent\begin{minipage}[t]{0.6\linewidth}\vspace{0pt}\par\noindent\hangindent6mm\hangafter1\makebox[6mm][l]{\textbf{1}}(a) Die bekannten Winkel zusammen: \ $85^\circ + 110^\circ + 72^\circ = 267^\circ$\par\smallskip\par\noindent\hangindent6mm\hangafter1\makebox[6mm][l]{\textbf{2}}$\delta = 360^\circ - 267^\circ = 93^\circ$\par\smallskip\par\noindent\hangindent6mm\hangafter1\makebox[6mm][l]{\textbf{3}}(b) $L$ und $N$ liegen neben der Achse $KM$: gleich groß. Für beide zusammen bleiben $360^\circ - 64^\circ - 100^\circ = 196^\circ$\par\smallskip\par\noindent\hangindent6mm\hangafter1\makebox[6mm][l]{\textbf{4}}je $196^\circ : 2 = 98^\circ$\par\smallskip\end{minipage}\hfill
\begin{minipage}[t]{0.37\linewidth}\vspace{0pt}\centering \begin{tikzpicture}[line width=0.7pt,font=\footnotesize]\draw (0.000,0.000) -- (2.400,0.000) -- (2.944,1.496) -- (0.139,1.594) -- cycle;\draw[mbgrau,line width=0.5pt] (0.380,0.000) arc (0.00:85.00:0.38);\node[inner sep=0pt] at (0.529,0.485) {$85^\circ$};\draw[mbgrau,line width=0.5pt] (2.530,0.357) arc (70.00:180.00:0.38);\node[inner sep=0pt] at (1.970,0.615) {$110^\circ$};\draw[mbgrau,line width=0.5pt] (2.565,1.509) arc (178.00:250.00:0.38);\node[inner sep=0pt] at (2.345,1.092) {$72^\circ$};\draw[mbgrau,line width=0.5pt] (0.106,1.215) arc (-95.00:-2.00:0.38);\node[inner sep=0pt] at (0.551,1.129) {$\delta$};\node[inner sep=0pt] at (-0.184,-0.169) {$A$};\node[inner sep=0pt] at (2.543,-0.205) {$B$};\node[inner sep=0pt] at (3.152,1.636) {$C$};\node[inner sep=0pt] at (-0.026,1.781) {$D$};\draw[line width=0.6pt] (5.310,1.633) -- (5.140,1.527);\draw[line width=0.6pt] (4.460,1.527) -- (4.290,1.633);\draw[line width=0.6pt] (5.134,0.597) -- (5.262,0.444);\draw[line width=0.6pt] (5.188,0.642) -- (5.316,0.489);\draw[line width=0.6pt] (4.338,0.444) -- (4.466,0.597);\draw[line width=0.6pt] (4.284,0.489) -- (4.412,0.642);\draw (4.800,2.260) -- (5.650,0.900) -- (4.800,0.187) -- (3.950,0.900) -- cycle;\draw[mbgrau,line width=0.5pt] (4.599,1.938) arc (-122.00:-58.00:0.38);\node[inner sep=0pt] at (4.800,1.635) {$64^\circ$};\draw[mbgrau,line width=0.5pt] (5.091,0.431) arc (40.00:140.00:0.38);\node[inner sep=0pt] at (4.800,0.782) {$100^\circ$};\node[inner sep=0pt] at (4.800,2.510) {$K$};\node[inner sep=0pt] at (5.897,0.861) {$L$};\node[inner sep=0pt] at (4.800,-0.063) {$M$};\node[inner sep=0pt] at (3.703,0.861) {$N$};\end{tikzpicture}\end{minipage}\par\medskip
\noindent\textbf{Aufgaben}\hspace{1em}{\small\color{mbgrau}\karohinweis{C}{Rechne unten im Karo.}{Rechne auf einem extra Blatt.} Vergleiche jede Aufgabe gleich mit dem Lösungsheft.}\par\smallskip
\aufgg{\hangindent7mm\hangafter1\nr{1.}Berechne den Winkel $\gamma$ im Viereck $ABCD$.}{\begin{tikzpicture}[line width=0.7pt,font=\footnotesize]\draw (0.000,0.000) -- (2.800,0.000) -- (3.586,1.123) -- (0.312,1.467) -- cycle;\draw[mbgrau,line width=0.5pt] (0.380,0.000) arc (0.00:78.00:0.38);\node[inner sep=0pt] at (0.560,0.454) {$78^\circ$};\draw[mbgrau,line width=0.5pt] (3.018,0.311) arc (55.00:180.00:0.38);\node[inner sep=0pt] at (2.466,0.643) {$125^\circ$};\draw[mbgrau,line width=0.5pt] (3.208,1.163) arc (174.00:235.00:0.38);\node[inner sep=0pt] at (3.042,0.875) {$\gamma$};\draw[mbgrau,line width=0.5pt] (0.233,1.096) arc (-102.00:-6.00:0.38);\node[inner sep=0pt] at (0.725,0.899) {$96^\circ$};\node[inner sep=0pt] at (-0.194,-0.157) {$A$};\node[inner sep=0pt] at (2.915,-0.222) {$B$};\node[inner sep=0pt] at (3.814,1.227) {$C$};\node[inner sep=0pt] at (0.165,1.669) {$D$};\end{tikzpicture}}{}{}
\medskip
\aufgg{\hangindent7mm\hangafter1\nr{2.}$ABCD$ ist ein Drachen mit der Symmetrieachse $AC$. Berechne $\beta$ und $\delta$.}{\begin{tikzpicture}[line width=0.7pt,font=\footnotesize]\draw[line width=0.6pt] (-1.084,-0.460) -- (-1.172,-0.640);\draw[line width=0.6pt] (-1.172,0.640) -- (-1.084,0.460);\draw[line width=0.6pt] (0.434,-0.572) -- (0.263,-0.468);\draw[line width=0.6pt] (0.398,-0.632) -- (0.227,-0.528);\draw[line width=0.6pt] (0.263,0.468) -- (0.434,0.572);\draw[line width=0.6pt] (0.227,0.528) -- (0.398,0.632);\draw (-2.255,0.000) -- (0.000,-1.100) -- (0.661,0.000) -- (0.000,1.100) -- cycle;\draw[mbgrau,line width=0.5pt] (-1.914,-0.167) arc (-26.00:26.00:0.38);\node[inner sep=0pt] at (-1.569,0.000) {$52^\circ$};\draw[mbgrau,line width=0.5pt] (0.196,-0.774) arc (59.00:154.00:0.38);\node[inner sep=0pt] at (-0.175,-0.510) {$\beta$};\draw[mbgrau,line width=0.5pt] (0.465,0.326) arc (121.00:239.00:0.38);\node[inner sep=0pt] at (-0.112,0.000) {$118^\circ$};\draw[mbgrau,line width=0.5pt] (-0.342,0.933) arc (-154.00:-59.00:0.38);\node[inner sep=0pt] at (-0.175,0.510) {$\delta$};\node[inner sep=0pt] at (-2.505,0.000) {$A$};\node[inner sep=0pt] at (0.071,-1.340) {$B$};\node[inner sep=0pt] at (0.911,0.000) {$C$};\node[inner sep=0pt] at (0.071,1.340) {$D$};\end{tikzpicture}}{}{}
\medskip
\aufg{\hangindent7mm\hangafter1\nr{3.}Welche Aussage gilt für jede Raute? Kreuze an.\par\smallskip\leftskip7mm \kreuz{Alle vier Seiten sind gleich lang.}\ \kreuz{Alle vier Winkel sind gleich groß.}\par\kreuz{Nur zwei Seiten sind parallel.}\ \kreuz{Benachbarte Winkel sind gleich groß.}}{}{}
\medskip
\aufgg{\hangindent7mm\hangafter1\nr{4.}Paul baut einen Flugdrachen aus zwei Leisten; die lange Leiste ist die Symmetrieachse. Im Bauplan steht: oben $84^\circ$, unten $60^\circ$, links und rechts je $106^\circ$. Kann das stimmen? Trag die richtigen Winkel links und rechts ins Bild ein.}{\begin{tikzpicture}[line width=0.7pt,font=\footnotesize]\fill[black!8] (0.000,1.111) -- (1.000,0.000) -- (0.000,-1.732) -- (-1.000,0.000) -- cycle;\draw (0.000,1.111) -- (1.000,0.000) -- (0.000,-1.732) -- (-1.000,0.000) -- cycle;\draw[line width=0.5pt] (0.000,-1.732) .. controls (0.350,-2.032) and (0.100,-2.382) .. (0.450,-2.682);\draw[mbgrau,line width=0.5pt] (-0.254,0.828) arc (-132.00:-48.00:0.38);\node[inner sep=0pt] at (0.000,0.516) {$84^\circ$};\draw[mbgrau,line width=0.5pt] (0.190,-1.403) arc (60.00:120.00:0.38);\node[inner sep=0pt] at (0.000,-1.077) {$60^\circ$};\draw[mbgrau,line width=0.5pt] (0.746,0.282) arc (132.00:240.00:0.38);\draw[mbgrau,line width=0.5pt] (-0.810,-0.329) arc (-60.00:48.00:0.38);\node[inner sep=0pt] at (1.420,0.000) {\textbf{?}};\node[inner sep=0pt] at (-1.420,0.000) {\textbf{?}};\end{tikzpicture}}{}{{\scriptsize\color{mbgrau}Ziel\hspace{2mm}}}
\medskip
\par\vspace{3mm}\setlength{\kr}{\dimexpr\pagegoal-\pagetotal-10mm\relax}%
\ifdim\kr>12mm\karomerk{k}{C}\noindent{\scriptsize\color{mbgrau}Platz zum Rechnen}\par\nointerlineskip\vspace{1mm}\noindent
\begin{tikzpicture}\clip (0,0) rectangle ({\linewidth-0.5mm},\kr);\draw[black!16,line width=0.3pt,step=5mm] (0,0) grid (\linewidth,\kr);\end{tikzpicture}\fi\par
\clearpage\label{s:D}\noindent\colorbox{black!12}{\makebox[10mm]{\Large\bfseries\strut D}}\hspace{3mm}{\Large\bfseries Winkel in Teildreiecken}\par\smallskip
\noindent Teilt eine Linie (Höhe, Diagonale) die Figur, rechnest du im Teildreieck, in dem der gesuchte Winkel liegt. Wo zwei Teildreiecke auf einer Geraden aneinanderstoßen, ergänzen sich ihre Winkel zu $180^\circ$.\par\medskip
\noindent\textbf{Formel}\par\nopagebreak\noindent\hspace*{4mm}\begin{minipage}{\dimexpr\linewidth-4mm\relax}Teildreieck: zusammen $180^\circ$ \qquad auf der Geraden: Nebenwinkel $= 180^\circ -$ Winkel\par {\small Höhe: Bei ihrem Fußpunkt sind beide Winkel $90^\circ$.}\end{minipage}\par\medskip
\noindent\textbf{Beispiel}\quad Im Dreieck $ABC$ liegt $D$ auf der Seite $AB$. Berechne den Winkel $\varepsilon$ zwischen $CD$ und $CB$.\par\smallskip
\noindent\begin{minipage}[t]{0.6\linewidth}\vspace{0pt}\par\noindent\hangindent6mm\hangafter1\makebox[6mm][l]{\textbf{1}}Teildreieck $ADC$: \ Winkel bei $D$ $= 180^\circ - 40^\circ - 35^\circ = 105^\circ$\par\smallskip\par\noindent\hangindent6mm\hangafter1\makebox[6mm][l]{\textbf{2}}$D$ liegt auf der Geraden $AB$: \textbf{Nebenwinkel}. Winkel $BDC = 180^\circ - 105^\circ = 75^\circ$\par\smallskip\par\noindent\hangindent6mm\hangafter1\makebox[6mm][l]{\textbf{3}}Teildreieck $DBC$, darin liegt $\varepsilon$: \ $\varepsilon = 180^\circ - 75^\circ - 62^\circ = 43^\circ$\par\smallskip\par\noindent\hangindent6mm\hangafter1\makebox[6mm][l]{\textbf{4}}Probe: $\gamma = 35^\circ + 43^\circ = 78^\circ$; \ $40^\circ + 62^\circ + 78^\circ = 180^\circ$\par\smallskip\end{minipage}\hfill
\begin{minipage}[t]{0.37\linewidth}\vspace{0pt}\centering \begin{tikzpicture}[line width=0.7pt,font=\footnotesize]\draw (0.000,0.000) -- (3.600,0.000) -- (2.489,2.089) -- cycle;\draw (2.489,2.089) -- (1.930,0.000);\draw[mbgrau,line width=0.5pt] (0.380,0.000) arc (0.00:40.00:0.38);\node[inner sep=0pt] at (0.788,0.287) {$40^\circ$};\draw[mbgrau,line width=0.5pt] (3.422,0.336) arc (118.00:180.00:0.38);\node[inner sep=0pt] at (2.980,0.373) {$62^\circ$};\draw[mbgrau,line width=0.5pt] (2.198,1.845) arc (-140.00:-105.00:0.38);\node[inner sep=0pt] at (1.922,1.199) {$35^\circ$};\draw[mbgrau,line width=0.5pt] (2.391,1.722) arc (-105.00:-62.00:0.38);\node[inner sep=0pt] at (2.558,1.488) {$\varepsilon$};\node[inner sep=0pt] at (-0.235,-0.086) {$A$};\node[inner sep=0pt] at (3.814,-0.129) {$B$};\node[inner sep=0pt] at (2.537,2.334) {$C$};\node[inner sep=0pt] at (1.930,-0.220) {$D$};\end{tikzpicture}\end{minipage}\par\medskip
\noindent\textbf{Aufgaben}\hspace{1em}{\small\color{mbgrau}\karohinweis{D}{Rechne unten im Karo.}{Rechne auf einem extra Blatt.} Vergleiche jede Aufgabe gleich mit dem Lösungsheft.}\par\smallskip
\aufgg{\hangindent7mm\hangafter1\nr{1.}$CD$ ist die Höhe im Dreieck $ABC$. Berechne $\varepsilon$, $\varphi$ und den ganzen Winkel $\gamma$ bei $C$.}{\begin{tikzpicture}[line width=0.7pt,font=\footnotesize]\draw (0.000,0.000) -- (3.243,0.000) -- (1.172,2.300) -- cycle;\draw (1.172,2.300) -- (1.172,0.000);\draw[mbgrau,line width=0.5pt] (0.380,0.000) arc (0.00:63.00:0.38);\node[inner sep=0pt] at (0.617,0.378) {$63^\circ$};\draw[mbgrau,line width=0.5pt] (2.989,0.282) arc (132.00:180.00:0.38);\node[inner sep=0pt] at (2.584,0.293) {$48^\circ$};\draw[mbgrau,line width=0.5pt] (1.172,0.300) arc (90.00:180.00:0.3);\fill (1.076,0.095) circle (0.8pt);\draw[mbgrau,line width=0.5pt] (0.999,1.961) arc (-117.00:-90.00:0.38);\node[inner sep=0pt] at (0.980,1.500) {$\varepsilon$};\draw[mbgrau,line width=0.5pt] (1.172,1.920) arc (-90.00:-48.00:0.38);\node[inner sep=0pt] at (1.393,1.723) {$\varphi$};\node[inner sep=0pt] at (-0.213,-0.131) {$A$};\node[inner sep=0pt] at (3.471,-0.102) {$B$};\node[inner sep=0pt] at (1.139,2.548) {$C$};\node[inner sep=0pt] at (1.172,-0.220) {$D$};\end{tikzpicture}}{}{}
\medskip
\aufgg{\hangindent7mm\hangafter1\nr{2.}$D$ liegt auf der Seite $AB$. Berechne den Winkel $\varepsilon$.}{\begin{tikzpicture}[line width=0.7pt,font=\footnotesize]\draw (0.000,0.000) -- (3.400,0.000) -- (2.679,2.093) -- cycle;\draw (2.679,2.093) -- (1.658,0.000);\draw[mbgrau,line width=0.5pt] (0.380,0.000) arc (0.00:38.00:0.38);\node[inner sep=0pt] at (0.849,0.292) {$38^\circ$};\draw[mbgrau,line width=0.5pt] (2.038,0.000) arc (0.00:64.00:0.38);\node[inner sep=0pt] at (2.272,0.383) {$64^\circ$};\draw[mbgrau,line width=0.5pt] (2.380,1.859) arc (-142.00:-116.00:0.38);\node[inner sep=0pt] at (2.118,1.400) {$\varepsilon$};\node[inner sep=0pt] at (-0.236,-0.081) {$A$};\node[inner sep=0pt] at (3.604,-0.145) {$B$};\node[inner sep=0pt] at (2.750,2.333) {$C$};\node[inner sep=0pt] at (1.658,-0.220) {$D$};\end{tikzpicture}}{}{}
\medskip
\aufgg{\hangindent7mm\hangafter1\nr{3.}$CD$ steht senkrecht auf $AB$, und $AD$, $DC$ und $DB$ sind gleich lang (Striche). Begründe, dass der Winkel $\gamma$ bei $C$ ein rechter ist.}{\begin{tikzpicture}[line width=0.7pt,font=\footnotesize]\draw (0.000,0.000) -- (3.000,0.000) -- (1.500,1.500) -- cycle;\draw (1.500,1.500) -- (1.500,0.000);\draw[line width=0.6pt] (0.750,0.100) -- (0.750,-0.100);\draw[line width=0.6pt] (2.250,0.100) -- (2.250,-0.100);\draw[line width=0.6pt] (1.400,0.750) -- (1.600,0.750);\draw[mbgrau,line width=0.5pt] (1.500,0.300) arc (90.00:180.00:0.3);\fill (1.405,0.095) circle (0.8pt);\node[inner sep=0pt] at (-0.231,-0.096) {$A$};\node[inner sep=0pt] at (3.231,-0.096) {$B$};\node[inner sep=0pt] at (1.500,1.750) {$C$};\node[inner sep=0pt] at (1.500,-0.220) {$D$};\end{tikzpicture}}{}{}
\medskip
\aufg{\hangindent7mm\hangafter1\nr{4.}In der Mitte eines Zelts steht eine senkrechte Stange, 1,4\,m hoch. Von ihrer Spitze gehen die beiden Zeltwände gerade zum Boden; sie sind links und rechts je 1,4\,m vom Fuß der Stange entfernt befestigt. Bilden die Zeltwände oben einen rechten Winkel? Begründe.}{}{{\scriptsize\color{mbgrau}Ziel\hspace{2mm}}}
\medskip
\par\vspace{3mm}\setlength{\kr}{\dimexpr\pagegoal-\pagetotal-10mm\relax}%
\ifdim\kr>12mm\karomerk{k}{D}\noindent{\scriptsize\color{mbgrau}Platz zum Rechnen}\par\nointerlineskip\vspace{1mm}\noindent
\begin{tikzpicture}\clip (0,0) rectangle ({\linewidth-0.5mm},\kr);\draw[black!16,line width=0.3pt,step=5mm] (0,0) grid (\linewidth,\kr);\end{tikzpicture}\fi\par
\clearpage\label{s:T}\noindent\colorbox{black!12}{\makebox[10mm]{\Large\bfseries\strut T}}\hspace{3mm}{\Large\bfseries Probetest}\par\smallskip
\noindent Gemischt wie in der Prüfung, ohne Beispiel. Die Skizzen sind nicht maßstabsgerecht: rechnen, nicht messen. Etwa 15 Minuten.\par\medskip
\noindent\textbf{Aufgaben}\hspace{1em}{\small\color{mbgrau}\karohinweis{T}{Rechne unten im Karo.}{Rechne auf einem extra Blatt.} Vergleiche jede Aufgabe gleich mit dem Lösungsheft.}\par\smallskip
\aufg{\hangindent7mm\hangafter1\nr{1.}Ein Dreieck hat die Winkel $46{,}5^\circ$ und $71{,}8^\circ$. Wie groß ist der dritte Winkel?\par\smallskip\leftskip7mm \kreuz{$61{,}7^\circ$}\ \kreuz{$62{,}7^\circ$}\par\kreuz{$118{,}3^\circ$}\ \kreuz{$51{,}7^\circ$}}{}{}
\medskip
\aufgg{\hangindent7mm\hangafter1\nr{2.}Im Trapez $ABCD$ ist $AB \parallel CD$. Berechne $\delta$.}{\begin{tikzpicture}[line width=0.7pt,font=\footnotesize]\draw (0.000,0.000) -- (3.200,0.000) -- (2.517,1.400) -- (0.482,1.400) -- cycle;\draw[-{Stealth[length=1.8mm,width=1.6mm]}] (1.600,0.000) -- (1.730,0.000);\draw[-{Stealth[length=1.8mm,width=1.6mm]}] (1.500,1.400) -- (1.630,1.400);\draw[mbgrau,line width=0.5pt] (0.380,0.000) arc (0.00:71.00:0.38);\node[inner sep=0pt] at (0.588,0.420) {$71^\circ$};\draw[mbgrau,line width=0.5pt] (3.033,0.342) arc (116.00:180.00:0.38);\node[inner sep=0pt] at (2.587,0.383) {$64^\circ$};\draw[mbgrau,line width=0.5pt] (2.137,1.400) arc (180.00:296.00:0.38);\node[inner sep=0pt] at (2.125,0.772) {$116^\circ$};\draw[mbgrau,line width=0.5pt] (0.358,1.041) arc (-109.00:0.00:0.38);\node[inner sep=0pt] at (0.843,0.893) {$\delta$};\node[inner sep=0pt] at (-0.204,-0.145) {$A$};\node[inner sep=0pt] at (3.412,-0.132) {$B$};\node[inner sep=0pt] at (2.650,1.612) {$C$};\node[inner sep=0pt] at (0.337,1.604) {$D$};\end{tikzpicture}}{}{}
\medskip
\aufgg{\hangindent7mm\hangafter1\nr{3.}Die Seite $AB$ ist über $B$ hinaus verlängert. Berechne $\gamma$.}{\begin{tikzpicture}[line width=0.7pt,font=\footnotesize]\draw (0.000,0.000) -- (2.600,0.000) -- (1.656,1.776) -- cycle;\draw (2.600,0.000) -- (3.600,0.000);\draw[mbgrau,line width=0.5pt] (0.380,0.000) arc (0.00:47.00:0.38);\node[inner sep=0pt] at (0.716,0.312) {$47^\circ$};\draw[mbgrau,line width=0.5pt] (2.980,0.000) arc (0.00:118.00:0.38);\node[inner sep=0pt] at (2.980,0.632) {$118^\circ$};\draw[mbgrau,line width=0.5pt] (1.397,1.498) arc (-133.00:-62.00:0.38);\node[inner sep=0pt] at (1.577,1.175) {$\gamma$};\node[inner sep=0pt] at (-0.229,-0.100) {$A$};\node[inner sep=0pt] at (2.891,-0.175) {$B$};\node[inner sep=0pt] at (1.688,2.024) {$C$};\end{tikzpicture}}{}{}
\medskip
\aufg{\hangindent7mm\hangafter1\nr{4.}Welche Aussage gilt für jeden Drachen? Kreuze an.\par\smallskip\leftskip7mm \kreuzl{Die Diagonalen stehen senkrecht aufeinander.}\kreuzl{Gegenüberliegende Seiten sind parallel.}\kreuzl{Alle vier Seiten sind gleich lang.}\kreuzl{Alle vier Winkel sind gleich groß.}}{}{}
\medskip
\aufg{\hangindent7mm\hangafter1\nr{5.}Zwei Häuser haben Satteldächer; links und rechts sind die Dachflächen gleich lang. Bei Haus 1 sind die Dachflächen um $35^\circ$ gegen den Dachboden geneigt. Bei Haus 2 misst der Winkel oben am First $104^\circ$. Welches Dach ist steiler? Um wie viel Grad?}{}{{\scriptsize\color{mbgrau}Ziel\hspace{2mm}}}
\medskip
\par\vspace{3mm}\setlength{\kr}{\dimexpr\pagegoal-\pagetotal-10mm\relax}%
\ifdim\kr>12mm\karomerk{k}{T}\noindent{\scriptsize\color{mbgrau}Platz zum Rechnen}\par\nointerlineskip\vspace{1mm}\noindent
\begin{tikzpicture}\clip (0,0) rectangle ({\linewidth-0.5mm},\kr);\draw[black!16,line width=0.3pt,step=5mm] (0,0) grid (\linewidth,\kr);\end{tikzpicture}\fi\par
\end{document}
